\subsection{The actual ordinary family and all arithmetic points}
\label{subsec:selmer-ordinary-families}

The original Hecke ideal satisfies the full family contract.  The
additional argument supplies classical control at every algebraic
weight, density at each fixed weight, the distinguished-prime condition,
and the exact decomposition and twist statements
\cite{SelmerFamiliesSource}.

\begin{theorem}\label{thm:selmer-ordinary-families}
Retain the original odd prime $\ell$, Hida setting $L/F$, level $V$
with $V.S'=S$, large coefficient DVR $O$, rank $n\ge1$, residual
datum, universal ordinary pair $B$, and weight ring
$A=O[[X_1,\ldots,X_D]]$, $D=n[F:\mathbb Q]$.  Retain also the
original type lattice, Hecke operators and residual summand.
Let $M$ be the constructed ordinary module.  For $M\ne0$ let
$T=A[H_i]\subseteq\operatorname{End}_A(M)$, let $f:B\twoheadrightarrow T$
be its constructed determinant map, and set $J=\ker f$; for $M=0$
set $J=B$.

This same $J$ satisfies every original \texttt{HidaProps} field:
\texttt{RingProps}, density at every fixed strictly increasing labeled
weight, PT, arithmetic-point purity and classicality, and
\texttt{TwistOk}.  With the proved P1/P3 propagation it satisfies the
original \texttt{Modular} contract and the already constructed
\texttt{HasAut}.  The \texttt{HidaU} application retains its hypothesis
that $L$ contains no primitive $\ell$th root of unity.
\end{theorem}

\begin{proof}
\emph{Full integral control.}
Fix an arithmetic weight specialization $x:A\to O_E$ and write
$\theta_x=\phi\chi_\lambda$, where $\phi$ has finite image and
$\lambda_{\tau,n-1-i}=h_{i,\tau}-i$ is dominant.  Use the previously
constructed saturated exterior-tensor lattice $L_\lambda$, its
primitive lowest vector $e_0$ and lowest-coordinate functional $e_0^*$.
The actual normalized operators
$R_a=\chi_\lambda(a)^{-1}\xi_\lambda(a)$ are integral, and $R_z$
contracts every nonlowest coordinate by the maximal ideal, for the
original $z=\operatorname{diag}(\varpi^{n-1},\ldots,1)$.

For $c\ge n$ beyond the compact-integrality depth and conductor of $\phi$,
let $C_c$ be the original classical coefficient lattice.  Untwisted
values at its finite double-coset representatives make it finite free.
The established arithmetic stabilizer calculation justifies arbitrary
values: a possible central stabilizer acts by $\theta_x$ times the
original type, hence trivially.  Evaluation is
\[
 \operatorname{ev}(F)(g)
   =e_0^*\xi_\lambda(g_\ell)^{-1}F(g),\qquad
 \operatorname{ev}(U_a^{\rm cl,norm}F)=U_a\operatorname{ev}(F).
                                                               \tag{OF1}
\]
It lies in the completed scalar module, with opposite-unipotent
invariance at each finite precision.  All original tame, spherical
and auxiliary actions commute with this map.

The original unnormalized coset sums give
\[
 U^mC_{c+e}\subseteq C_c\quad(m\ge e).                  \tag{OF2}
\]
Here the original $T_c$ and $\overline N(c)$ have depth $c+1$.
Translation by $N$ reindexes $N/(zNz^{-1})$, leaving an $N$-factor
after $z$; conjugation by $T$ permutes that set and gives precisely
the factor $\phi(t)^{-1}$.  For $y\in\overline N(c)$ and a coset
representative $u$, integral Gaussian elimination gives
$u^{-1}yu=n_+tn_-$, with $n_+\in N\cap K(c+1)$,
$t\in T_c$, and $n_-\in\overline N(c)$.
Since $c\ge n$,
\[
 y u z=u z\,(z^{-1}n_+z)t(z^{-1}n_-z),\qquad
 z^{-1}n_+z\in N,\quad z^{-1}n_-z\in\overline N(c+1).
\]
The positive-root exponents are at most $n-1$ and the negative ones
increase depth by at least one.  Thus every summand is fixed by $y$
on $C_{c+1}$, because $\phi(t)=1$.  This proves one-step lowering,
and iteration proves (OF2).  The original normalization
$\chi_\lambda(z)^{-1}$ is one common scalar; it introduces no new
torus factor or denominator in this calculation.
In untwisted coordinates
each summand is $\xi_\lambda(u)R_z$ times a value.  Therefore
\[
 F\in\varpi_E^{a-1}C,\quad
 \operatorname{ev}(F)\in\varpi_E^a\text{ functions}
 \quad\Longrightarrow\quad UF\in\varpi_E^aC.           \tag{OF3}
\]
Here $\xi_\lambda(u)$ is integral, the lowest coordinate already
has the required divisibility, and every other column contracts.
The same statements hold for the product over the $\ell$ places.

Let $e_c$ be the integral factorial-limit ordinary projector on $C_c$.
Both $U$ and $U^{-1}$ preserve $C_c^{\rm ord}=e_cC_c$.
Evaluation maps this lattice into $M_x=M\otimes_A O_E$.  It is
injective: applying (OF3) repeatedly to $U^{-r}F$ for a kernel vector
puts $F$ in every $\varpi_E^rC_c$.

For surjectivity take $f_0\in M_x$.  Choose $c'\ge c$ so that
$f_0$ and $\xi_\lambda$ are invariant modulo $\varpi_E$ under
$\overline N(c')$.  Prescribe the untwisted values $f_0(g)e_0$
at the original level-$c'$ representatives, in every type coordinate.
The resulting $F'\in C_{c'}$ has
$\operatorname{ev}(F')\equiv f_0\pmod{\varpi_E}$ everywhere:
the exact $N,T$ and tame transformation rules, and the chosen two
$\overline N(c')$ congruences, propagate the equality from the
representatives.  Choose $r\ge c'-c$ divisible by the period of
$f_0$ under $U$ modulo $\varpi_E$.  Then $G=U^rF'\in C_c$ and
$\operatorname{ev}(G)\equiv f_0$.  Applying $e_c$ preserves this
congruence.  The finite image lattice $L=\operatorname{ev}(C_c^{\rm ord})$
thus satisfies $L+\varpi_E M_x=M_x$; Nakayama gives
\[
       C_c^{\rm ord}\xrightarrow[\operatorname{ev}]{\ \sim\ }
                     M\otimes_A O_E\qquad(c\ge c_0).             \tag{OF4}
\]
This is the full specialization interface, with the original integral
coefficient lattice and all Hecke and residual actions.  Neither
the depth $c_0$ nor the coefficient rank is declared uniformly bounded.

\emph{Density at each weight, and occurrence of every arithmetic point.}
Fix $h$.  Local reciprocity describes its algebraic unit character
as $u\mapsto\prod_\tau\tau(u)^{-h_{i,\tau}}$.  Adjust its finite
torsion restriction and residual character by a finite character,
using the finite-torsion plus free-pro-$\ell$ decomposition of local
units.  Vary the $\ell$-power finite characters of the free factor.
After a fixed finite coefficient extension the original coordinates are
\[
          1+X_j=u_j\zeta_j,\qquad
 u_j\in1+\mathfrak m,\quad\zeta_j\in\mu_{\ell^\infty}.           \tag{OF5}
\]
Translation by $u_j$ is a power-series automorphism.  A nonzero
one-variable series, after removing its common coefficient valuation,
has its first unit-coefficient term uniquely minimal at roots
$\zeta_{\ell^a}-1$ of sufficiently small positive valuation.
Induction gives multivariable density, still allowing finitely many
excluded roots at each step.  Exclude equalities between the finite
unit characters at one free generator for each place.  These are
exactly good points of the fixed $h$.

For a nonzero $t\in T$ choose $v$ with $tv\ne0$ and such a point
where a coordinate survives.  By (OF4) this is a classical space.
Its original ordinary summands inject into full Jacquet coinvariants;
distinct finite unit characters select one Jacquet line.  The
principal-series and geometric-lemma proof already used for the
Hecke determinant therefore makes every original normalized ordinary
operator scalar on each automorphic summand.  Spherical operators
are scalar too; all global multiplicities remain.  A classical
eigencharacter detects $t$.  This proves the literal fixed-weight
\texttt{Dense} condition.

Now let $x:B\to\overline{\mathbb Q}_\ell$ kill $J$ and be arithmetic,
without a good condition.  It induces $y:T\to\overline{\mathbb Q}_\ell$.
The established automatic-continuity argument puts its image in a
finite $O_E$.  In particular every coefficient is a residue lift from
$O$ plus a maximal-ideal element, so its image is integral before the
uniform valuation bound from finitely many ideal generators is used.
The finite faithful $T$-module $M$ has full support; localization and
Nakayama give a nonzero $y$-summand of $M\otimes_A E$.
Its local Artinian operator algebra has nonzero socle, yielding a
common eigenvector.  Thus no faithfulness assertion after a nonflat
weight change is needed.  Equation (OF4) realizes this character in
an actual classical automorphic constituent $\sigma$ on the original
unitary group, even when ramified operators are not semisimple.

\emph{All arithmetic-point decompositions.}
Labesse's Corollary 5.3 applies: $[F:\mathbb Q]\ge2$ since the CM
extension is unramified at every finite place, and compact infinity
supplies discrete series and condition $(*)$.
It gives discrete conjugate-self-dual base-change blocks, with matching
infinitesimal character and split/unramified parameters
\cite[Corollary 5.3]{SelmerFamiliesLabesse}.
Write these as $\operatorname{Speh}(\tau_a,m_a)$ by
M\oe glin--Waldspurger, where $\tau_a$ is unitary cuspidal of rank $d_a$,
$m_a\ge1$, and $\sum_a d_am_a=n$
\cite[pp.~605--607]{SelmerFamiliesMW}.  Uniqueness of cuspidal support
gives $\tau_a^c=\tau_a^\vee$.  Set
\[
 \pi_a=\tau_a|\det|^{(d_a+m_a-n-1)/2},\qquad
 \operatorname{algNorm}_{d_a}(\pi_a)
                 =\tau_a|\det|^{(m_a-n)/2}.             \tag{OF6}
\]
At the top of a Speh string the entry of $\tau_a$ is
$(n-m_a)/2-h_i$; the shift in (OF6) makes it $(d_a-1)/2-h_i$.
Thus $\pi_a$ is regular algebraic in precisely rank $d_a$.
It is polarized with character
$\delta_{L/F}^{d_a}|\cdot|^{d_a+m_a-n-1}$, retaining the real sign.
BLGGT's adopted Galois theorem gives its continuous semisimple $r_a$
\cite[Theorem 2.1.1]{SelmerFamiliesBLGGT}.  The exact representation
at $x$ is
\[
       \rho_x=\bigoplus_a\,
                   \bigoplus_{j=0}^{m_a-1}r_a\epsilon^{-j}.       \tag{OF7}
\]
The Speh inducing exponents $(m_a-1)/2-j$ and (OF6) give exactly
the rank-$n$ Frobenius polynomial.  Unramified comparison outside
$S$ and Chebotarev identify its full determinant with $xD_B$.
Every rank and multiplicity in the source decomposition is retained.

If all $m_a=1$, the multiplier of $r_a$ is $\epsilon^{1-n}$.
BLGGT purity then has weight $n-1$: a pure weight $w$ satisfies
$w=2(n-1)-w$ under this polarization.  The source \texttt{FrobEig}
predicate includes unramifiedness, so at every such place away from
$\ell$ all eigenvalues have absolute value $(Nw)^{(n-1)/2}$.
This proves \texttt{arithmeticPointPurity}.  If a representation with
determinant $xD_B$ is absolutely irreducible, character theory identifies
it with (OF7); there is one block, $m_1=1$, $d_1=n$.
Then $\pi_1=\tau_1$ without a half-norm shift, which gives the exact
regular algebraic conjugate-self-dual cuspidality statement.  Its
fixed conjugation to (OF7) also supplies continuity if that was not
an input.

\emph{The distinguished prime and the remaining ring fields.}
Here is the local point needed by \texttt{RingProps.unr0}.
Suppose $p\ne\ell$, $q\not\equiv1\pmod\ell$, and a continuous
ell-adic representation has unramified residual semisimplification
with Frobenius polynomial $(X-a)^n$.  In any stable lattice its
residual inertia image is an $\ell$-group: every element is unipotent,
and Cauchy's theorem applies.  Wild and prime-to-$\ell$ tame inertia
therefore vanish, including in all higher congruence quotients.
Write $T$ for the remaining pro-$\ell$ tame generator and $\Phi$
for arithmetic Frobenius, so $\Phi T\Phi^{-1}=T^q$.
The conjugation of $\overline\Phi$, a scalar times a unipotent matrix,
has $\ell$-power order on $\langle\overline T\rangle$.
If this cyclic group were nontrivial, raising to $q$ would have
nontrivial prime-to-$\ell$ order.  Thus $\overline T=1$.
The successive congruence equation is
\[
  (\operatorname{Ad}(\overline\Phi)-q)
                (T-1)/\varpi_E^r=0.                              \tag{OF8}
\]
Its operator is invertible because $\operatorname{Ad}(\overline\Phi)$
is unipotent.  Induction gives $T=1$.  No semisimple residual lattice
or rank restriction was imposed.

These are exactly the original \texttt{V0ok} conditions.  Apply this
fact at the separating weight-zero Galois points, then use coefficient
separation to obtain both literal determinant-kernel relations at
$I_0$, including the polarization star.  The other ring fields follow
from the existing determinant theorem: finiteness and torsion-freeness
over $A$, polarization, ordered ordinary relations and the prescribed
types.  Thus the entire \texttt{RingProps} contract holds.

\emph{Twist on the original unitary group.}
Let $\beta$ have the exact \texttt{TwistOk} hypotheses: finite
$\ell$-power order, $\beta^c=\beta^{-1}$, and triviality on all old
inertia and $I_0$.  Use the source convention
$v(\operatorname{art}g)=-\deg g$: arithmetic Frobenius maps to an
inverse uniformizer, whereas $\operatorname{art}(g_\varpi)=\varpi$
has degree $-1$.  Define
$b_w(a)=\iota\beta(\operatorname{rec}_{\rm geom}(a))$, using the
inverse of the arithmetic reciprocity map
\cite[V.5.2--5.3]{SelmerFamiliesCFT}.
The finite-character avatar identity has no algebraic correction,
so $b_w(\varpi)=\iota\beta(g_\varpi)$ exactly.
On $\mathbb A_F^\times$ it satisfies $b^2=1$, hence is
trivial because its order is odd.  If $2r\equiv1\pmod{\operatorname{ord}b}$,
define on the norm-one torus
\[
                 b_U(t)=b(t)^r.                                  \tag{OF9}
\]
It is trivial on $U(1)(F)$.  Equivalently $b_U(a/a^c)=b(a)$.
At a split place $b_U(u,u^{-1})=b_w(u)$, so this is the correct
character, without an erroneous square.

Multiply an actual classical eigenform for $x$ by $b_U(\det g)$.
This preserves the original rational group, infinity coefficient,
tame type and principal level at $v_0$.  Outside $S$ it preserves
hyperspecial level; at nonsplit unramified places $b_U$ is trivial
on the compact norm-one group.  At $\ell$ the finite torus character
changes by $b_w^{-1}$ and normalized $U_j$ is multiplied by
$\beta(g_\varpi)^j$, a unit.  The residual block is unchanged since
these $\ell$-power roots reduce to one.  Full control yields a
character of the original $T$.  Its determinant and its ordered
local characters are exactly the prescribed twists.  Dense Frobenius,
local reciprocity and topological generation by the full determinant
coefficients and ordered characters identify the induced point with
the given $x'$.  Hence $J\subseteq\ker x'$, proving \texttt{TwistOk}.

Finally PT is already proved in its complete original range.
Together the fields above give \texttt{HidaProps}; the established
arithmetic propagation supplies P1/P3 for the same ideal and actual
modules.  If $M=0$, all conditions requiring a point of $V(J)$ are
vacuous and the ring quotient is zero; the separate PT proof excludes
its input by constructing a nonzero ordinary vector.  This covers
the original zero branch as well.
\end{proof}

For an actual absolutely irreducible seed there is now a starting
component of dimension $1+n[F:\mathbb Q]$.  Indeed $T$ is finite and
$A$-torsion-free, so its minimal primes contract to zero: a nonzero
element of $A$ in a minimal associated prime would give torsion.
Each component therefore has the dimension of $A$.  Choose one below
the seed's kernel.  Stability follows separately from the seed's
nonzero $n^2$ trace-Gram discriminant, which stays outside every
smaller prime.  The remaining frontier is the exact normalized-Space
and \texttt{SeedIn} transfer followed by the global component argument
to \texttt{FinA}; no final uniform integral torsion bound is claimed.
